\documentclass[institutional,screen,none]{../cls/ua-engineering-v6-beamer}
% !TeX encoding = UTF-8
\UASetUniversity{Universidad Austral}
\UASetFaculty{Facultad de Ingeniería}
\UASetDepartment{Departamento de Computación}
\UASetCampus{Pilar}
\UASetCareer{Ingeniería Informática}
\UASetCommission{Comisión A}
\UASetProfessor{Equipo docente}
\UASetSemester{Segundo cuatrimestre 2026}
\UASetCourse{Sistemas Distribuidos}
\UASetDate{2026}


\UASetClassNumber{15}
\UASetClassTitle{Diseño integral de sistemas distribuidos}
\title{\UACourse}
\subtitle{Clase 15 --- Defensa arquitectónica, trade-offs y revisión de casos}

\begin{document}

{
\setbeamertemplate{footline}{}
\begin{frame}[plain]
\titlepage
\end{frame}
}

\UASectionSlide{Del mecanismo al sistema}

\begin{frame}{El cambio de nivel}
Durante el curso vimos:
\begin{itemize}
\item RPC, retries, idempotencia
\item consistencia, replicación, consenso
\item transacciones, logs, streaming
\item observabilidad, resiliencia, incidentes
\end{itemize}

\pause

\begin{block}{Pregunta}
¿Cómo integramos todo eso en un sistema coherente y defendible?
\end{block}
\end{frame}

\begin{frame}{Diseñar no es listar tecnologías}
Una propuesta arquitectónica no es:
\begin{itemize}
\item un stack de herramientas
\item un diagrama con cajas
\item una colección de patrones “modernos”
\end{itemize}

\pause

Es:
\begin{itemize}
\item una respuesta a un problema concreto
\item bajo fallas concretas
\item con garantías explícitas
\item y costos aceptados
\end{itemize}
\end{frame}

\UASectionSlide{Marco de defensa}

\begin{frame}{Marco del curso}
\[
D = (P, F, G, S, T)
\]

\pause

\begin{itemize}
\item $P$: problema
\item $F$: fallas asumidas
\item $G$: garantías buscadas
\item $S$: solución elegida
\item $T$: trade-offs aceptados
\end{itemize}
\end{frame}

\begin{frame}{Qué debe defenderse}
Toda propuesta seria debe justificar:
\begin{itemize}
\item límites del sistema
\item modelo de datos
\item consistencia requerida
\item estrategia de comunicación
\item estrategia de resiliencia
\item estrategia operacional
\end{itemize}
\end{frame}

\UASectionSlide{Trade-offs reales}

\begin{frame}{No existe arquitectura gratis}
Cada decisión gana algo y pierde algo.

\pause

Ejemplos:
\begin{itemize}
\item consistencia fuerte $\rightarrow$ más coordinación
\item microservicios $\rightarrow$ más autonomía, más complejidad distribuida
\item event-driven $\rightarrow$ más desacoplamiento, menos trazabilidad inmediata
\item resiliencia fuerte $\rightarrow$ más complejidad operativa
\end{itemize}
\end{frame}

\begin{frame}{Error típico}
\begin{alertblock}{}
Defender una arquitectura como si maximizara simultáneamente simplicidad, escalabilidad, consistencia, velocidad y costo bajo.
\end{alertblock}

\pause

Eso no es defensa técnica: es marketing.
\end{frame}

\UASectionSlide{Estructura de una defensa}

\begin{frame}{Secuencia razonable}
\begin{enumerate}
\item problema y contexto
\item restricciones y supuestos
\item requisitos no funcionales
\item diseño propuesto
\item fallas esperadas
\item trade-offs aceptados
\item observabilidad y operación
\end{enumerate}
\end{frame}

\begin{frame}{Qué valora un jurado técnico}
\begin{itemize}
\item claridad
\item coherencia interna
\item honestidad sobre límites
\item capacidad de priorizar
\item justificación de trade-offs
\end{itemize}
\end{frame}

\UASectionSlide{Revisión de casos}

\begin{frame}{Caso 1: e-commerce global}
Preguntas clave:
\begin{itemize}
\item ¿qué partes requieren consistencia fuerte?
\item ¿qué partes toleran eventual consistency?
\item ¿dónde usar eventos?
\item ¿cómo resolver checkout?
\end{itemize}
\end{frame}

\begin{frame}{Caso 2: plataforma financiera}
Preguntas clave:
\begin{itemize}
\item ¿qué operaciones exigen serialización fuerte?
\item ¿qué errores son inaceptables?
\item ¿qué latencia puede tolerarse?
\item ¿qué auditoría se exige?
\end{itemize}
\end{frame}

\begin{frame}{Caso 3: plataforma de streaming}
Preguntas clave:
\begin{itemize}
\item ¿qué volumen de eventos entra?
\item ¿qué semántica de entrega importa?
\item ¿qué parte es batch y qué parte es tiempo real?
\item ¿qué SLOs tienen sentido?
\end{itemize}
\end{frame}

\UASectionSlide{Objeciones típicas}

\begin{frame}{Preguntas hostiles esperables}
\begin{itemize}
\item ¿por qué no un monolito?
\item ¿por qué no consistencia fuerte en todo?
\item ¿qué pasa bajo partición?
\item ¿cómo observás este flujo?
\item ¿qué ocurre si falla este servicio crítico?
\item ¿cómo evitás duplicación?
\end{itemize}
\end{frame}

\begin{frame}{Buena defensa}
No responde con slogans.

\pause

Responde con:
\begin{itemize}
\item supuestos explícitos
\item garantías claras
\item límites reconocidos
\item elección consciente de costos
\end{itemize}
\end{frame}

\UASectionSlide{Cierre}

\begin{frame}{Idea central}
\begin{block}{}
Diseñar bien un sistema distribuido es poder defender por qué ese sistema, con esos costos y no otros, es razonable para ese dominio.
\end{block}
\end{frame}

\end{document}
